% Part: first-order-logic
% Chapter: completeness
% Section: lindenbaums-lemma

\documentclass[../../../include/open-logic-section]{subfiles}

\begin{document}

\iftag{FOL}
      {\olfileid{fol}{com}{lin}}
      {\olfileid{pl}{com}{lin}}

\olsection{لمّة ليندنبوم}

\begin{explain}
نبرهن الآن لمّة تبين أن كل مجموعة متسقة من ال!!{sentence}s محتواة في
مجموعة من الجمل ليست متسقة فحسب، بل !!{complete} أيضًا. ويعمل البرهان
بإضافة !!{sentence} واحدة في كل مرة، مع ضمان بقاء المجموعة متسقة في
كل خطوة. ونفعل ذلك بحيث تُضاف، لكل $!A$، إما $!A$ وإما $\lnot !A$ في
مرحلة ما. وعندئذ يحتوي اتحاد جميع مراحل هذا الإنشاء إما $!A$ وإما
نفيها~$\lnot !A$، فيكون كاملًا. وهو متسق أيضًا، لأننا نحرص في كل مرحلة
على ألا ندخل تناقضًا.
\end{explain}

\begin{lem}[لمّة ليندنبوم]
\ollabel{lem:lindenbaum} يمكن توسيع كل مجموعة متسقة~$\Gamma$ في
لغة~$\Lang{L}$ إلى مجموعة !!a{complete} ومتسقة~$\Gamma^*$.
\end{lem}

\begin{proof}
لتكن $\Gamma$ متسقة. وليكن $!A_0$ و$!A_1$ و\dots{} تعدادًا لجميع
ال!!{sentence}s في~$\Lang L$. عرّف $\Gamma_0 = \Gamma$، و
\[
\Gamma_{n+1} =
\begin{cases}
\Gamma_n \cup \{ !A_n \} & \textrm{إذا كانت $\Gamma_n \cup \{!A_n\}$
  متسقة؛} \\
\Gamma_n \cup \{ \lnot !A_n \} & \textrm{وخلاف ذلك.}
\end{cases}
\]
ولتكن $\Gamma^* = \bigcup_{n \geq 0} \Gamma_n$.

كل $\Gamma_n$ متسقة: فـ$\Gamma_0$ متسقة بالتعريف. وإذا كان
$\Gamma_{n+1} = \Gamma_n \cup \{!A_n\}$، فلأن المجموعة الأخيرة متسقة.
وإذا لم تكن متسقة، كان $\Gamma_{n+1} = \Gamma_n \cup \{\lnot !A_n\}$.
وعلينا أن نتحقق من اتساق $\Gamma_n \cup \{\lnot !A_n\}$. افترض أنها
ليست متسقة. عندئذ تكون \emph{كل من} $\Gamma_n \cup \{!A_n\}$ و
$\Gamma_n \cup \{\lnot !A_n\}$ غير متسقة. وهذا يعني أن $\Gamma_n$
غير متسقة بحسب
\iftag{FOL}{%
  \tagrefs{prfAX/{fol:axd:prv:prop:provability-exhaustive},
    prfSC/{fol:seq:prv:prop:provability-exhaustive},
    prfND/{fol:ntd:prv:prop:provability-exhaustive},
    prfTab/{fol:tab:prv:prop:provability-exhaustive}}}{%
  \tagrefs{prfAX/{pl:axd:prv:prop:provability-exhaustive},
    prfSC/{pl:seq:prv:prop:provability-exhaustive},
    prfND/{pl:ntd:prv:prop:provability-exhaustive},
    prfTab/{pl:tab:prv:prop:provability-exhaustive}}},
خلافًا لفرض الاستقراء.

لكل~$n$ ولكل $i < n$، لدينا $\Gamma_i \subseteq \Gamma_n$. وينتج ذلك
باستقراء بسيط في~$n$. فعندما $n=0$ لا يوجد $i < 0$، فتصح الدعوى
تلقائيًا. وفي خطوة الاستقراء، افترض صحتها لـ~$n$. ونبين أنه إذا كان
$i < n+1$، فإن $\Gamma_i \subseteq \Gamma_{n+1}$. بحسب الإنشاء، لدينا
$\Gamma_{n+1} = \Gamma_n \cup \{!A_n\}$ أو $= \Gamma_n \cup \{\lnot
!A_n\}$. ولذلك $\Gamma_n \subseteq \Gamma_{n+1}$. وإذا كان $i < n+1$،
فإن $\Gamma_i \subseteq \Gamma_n$ بحسب فرض الاستقراء (إذا كان
$i < n$)، أو بحسب الحقيقة البديهية $\Gamma_n \subseteq \Gamma_n$
(إذا كان $i = n$). ومن تعدية~$\subseteq$ نحصل على
$\Gamma_i \subseteq \Gamma_{n+1}$.

وينتج من ذلك أن $\Gamma^*$ متسقة. وبيان ذلك كما يأتي: لتكن
$\Gamma' \subseteq \Gamma^*$ منتهية. إذا كان $\Gamma'=\emptyset$، فهي
متسقة بداهة. وخلاف ذلك، فإن كل $!B \in \Gamma'$ تنتمي أيضًا
إلى~$\Gamma_i$ لبعض~$i$. وليكن $n$ أكبر هذه الفهارس. ولما كان
$\Gamma_i \subseteq \Gamma_n$ إذا كان $i \le n$، فإن كل
$!B \in \Gamma'$ تحقق $!B \in \Gamma_n$؛ أي إن
$\Gamma' \subseteq \Gamma_n$، و$\Gamma_n$~متسقة. ومن ثم تكون كل مجموعة
جزئية منتهية $\Gamma' \subseteq \Gamma^*$ متسقة. وبحسب \iftag{FOL}{%
  \tagrefs{prfAX/{fol:axd:ptn:prop:proves-compact},
    prfSC/{fol:seq:ptn:prop:proves-compact},
    prfND/{fol:ntd:ptn:prop:proves-compact},
    prfTab/{fol:tab:ptn:prop:proves-compact}}}{%
  \tagrefs{prfAX/{pl:axd:ptn:prop:proves-compact},
    prfSC/{pl:seq:ptn:prop:proves-compact},
    prfND/{pl:ntd:ptn:prop:proves-compact},
    prfTab/{pl:tab:ptn:prop:proves-compact}}}، تكون $\Gamma^*$ متسقة.

تظهر كل !!{sentence} من $\Frm[L]$ في القائمة المستعملة لتعريف
~$\Gamma^*$. وإذا كان $!A_n \notin \Gamma^*$، فلأن
$\Gamma_n \cup \{!A_n\}$ كانت غير متسقة. ولكن عندئذ يكون
$\lnot !A_n \in \Gamma^*$، ومن ثم تكون $\Gamma^*$ !!{complete}.
\end{proof}

\end{document}
